% Figure: returned optical power and detector signal-to-noise ratio against
% range for a pulsed time-of-flight LIDAR. The return power falls as 1/R^2 in
% the beam-overfills-target (point-target) regime, and the SNR falls with it
% until it crosses the detection threshold that fixes the maximum range.
\begin{figure}[htbp]
\centering
\begin{tikzpicture}
\begin{semilogyaxis}[
  width=13cm, height=7.2cm,
  xlabel={range $R$ (m)},
  ylabel={SNR (linear)},
  xmin=0, xmax=320, ymin=1, ymax=4000,
  axis lines=left,
  xtick={0,50,100,150,200,250,300},
  ytick={1,10,100,1000},
  tick label style={font=\scriptsize},
  label style={font=\small},
  legend style={font=\scriptsize, at={(0.98,0.97)}, anchor=north east, draw=none},
  enlarge x limits=false,
]
% point-target return SNR ~ K / R^2, normalised so SNR = 2500 at R = 20 m.
% K = 2500 * 20^2 = 1.0e6. Clamp domain away from R=0.
\addplot[engblue, very thick, samples=120, domain=20:320] {1.0e6/(x^2)};
\addlegendentry{return SNR ($\propto 1/R^2$)}
% detection threshold SNR = 6 (Q = 6 for BER-equivalent reliable ranging)
\addplot[engred, thick, dashed, domain=0:320] {6};
\addlegendentry{detection threshold (SNR $=6$)}
% maximum range where SNR = 6: R = sqrt(1.0e6/6) = 408 m, beyond plot; mark the
% in-plot working range crossing at SNR = 6 -> off scale, so mark R = 300.
\draw[engorange, dotted] (axis cs:300,1) -- (axis cs:300,12);
\node[engorange, font=\scriptsize, anchor=west] at (axis cs:240,30)
  {working range};
\node[enggray, font=\scriptsize, anchor=south west] at (axis cs:60,300)
  {$1/R^2$ point-target falloff};
\end{semilogyaxis}
\end{tikzpicture}
\caption[LIDAR return SNR against range]{Detector signal-to-noise ratio of a
pulsed time-of-flight LIDAR against range for a small (point-like) target that
the beam overfills, where the collected return power falls as $1/R^2$ and the
shot-noise-limited SNR falls with it. The horizontal line is the detection
threshold below which a return cannot be reliably timed; the range at which the
falling SNR curve crosses the threshold sets the maximum unambiguous range. A
target that fills the beam (a wall, the ground at grazing incidence) returns
power that falls only as $1/R^2$ through the receiver solid angle alone and
reaches farther, which is why a flat wall ranges farther than a wire at the
same distance.}
\label{fig:vol04ch11:lidar-range-snr}
\end{figure}
